\documentclass[11pt,a4paper]{article}

\usepackage[utf8]{inputenc}
\usepackage[T1]{fontenc}
\usepackage{lmodern}
\usepackage[margin=1in]{geometry}
\usepackage{hyperref}
\usepackage{xcolor}
\usepackage{listings}
\usepackage{enumitem}
\usepackage{graphicx}
\usepackage{titlesec}
\usepackage{tcolorbox}
\usepackage{array}
\usepackage{booktabs}
\usepackage{longtable}

% ---------- Hyperref ----------
\hypersetup{
    colorlinks=true,
    linkcolor=blue!60!black,
    urlcolor=blue!60!black,
    citecolor=blue!60!black,
    pdftitle={Deploying ML Models on AWS ECS and EKS},
    pdfauthor={Prabal}
}

% ---------- Listings ----------
\definecolor{codebg}{RGB}{245,245,245}
\definecolor{codekw}{RGB}{0,0,180}
\definecolor{codestr}{RGB}{163,21,21}
\definecolor{codecom}{RGB}{0,128,0}

\lstdefinestyle{pycode}{
    backgroundcolor=\color{codebg},
    basicstyle=\ttfamily\footnotesize,
    keywordstyle=\color{codekw}\bfseries,
    stringstyle=\color{codestr},
    commentstyle=\color{codecom}\itshape,
    language=Python, breaklines=true, showstringspaces=false,
    numbers=left, numberstyle=\tiny\color{gray},
    frame=single, rulecolor=\color{gray!30}, tabsize=2
}
\lstdefinestyle{bashcode}{
    backgroundcolor=\color{codebg},
    basicstyle=\ttfamily\footnotesize,
    keywordstyle=\color{codekw}\bfseries,
    stringstyle=\color{codestr},
    commentstyle=\color{codecom}\itshape,
    language=bash, breaklines=true, showstringspaces=false,
    frame=single, rulecolor=\color{gray!30}, tabsize=2
}
\lstdefinelanguage{docker}{
    keywords={FROM, RUN, CMD, ENTRYPOINT, WORKDIR, COPY, ADD, ENV, ARG, EXPOSE, VOLUME, USER, LABEL, HEALTHCHECK, SHELL},
    keywordstyle=\color{codekw}\bfseries, sensitive=true,
    comment=[l]{\#}, morestring=[b]",
}
\lstdefinestyle{dockercode}{
    backgroundcolor=\color{codebg}, basicstyle=\ttfamily\footnotesize,
    keywordstyle=\color{codekw}\bfseries, stringstyle=\color{codestr},
    commentstyle=\color{codecom}\itshape,
    language=docker, breaklines=true, showstringspaces=false,
    frame=single, rulecolor=\color{gray!30}, tabsize=2
}
\lstdefinelanguage{yaml}{
    keywords={apiVersion, kind, metadata, spec, name, namespace, replicas, selector, matchLabels, labels, template, containers, image, ports, containerPort, resources, requests, limits, memory, cpu, env, valueFrom, secretKeyRef, configMapKeyRef, livenessProbe, readinessProbe, httpGet, path, port, initialDelaySeconds, periodSeconds, type, clusterIP, targetPort, protocol, ingress, rules, host, http, paths, backend, service, number, strategy, rollingUpdate, maxSurge, maxUnavailable, nodeSelector, tolerations, affinity, serviceAccountName},
    keywordstyle=\color{codekw}\bfseries, sensitive=true,
    comment=[l]{\#}, morestring=[b]",
}
\lstdefinestyle{yamlcode}{
    backgroundcolor=\color{codebg}, basicstyle=\ttfamily\footnotesize,
    keywordstyle=\color{codekw}\bfseries, stringstyle=\color{codestr},
    commentstyle=\color{codecom}\itshape,
    language=yaml, breaklines=true, showstringspaces=false,
    frame=single, rulecolor=\color{gray!30}, tabsize=2
}
\lstdefinelanguage{json}{
    string=[s]{"}{"}, stringstyle=\color{codestr},
    comment=[l]{//}, commentstyle=\color{codecom},
}
\lstdefinestyle{jsoncode}{
    backgroundcolor=\color{codebg}, basicstyle=\ttfamily\footnotesize,
    language=json, breaklines=true, showstringspaces=false,
    frame=single, rulecolor=\color{gray!30}, tabsize=2
}

% ---------- Note boxes ----------
\newtcolorbox{notebox}{
    colback=blue!5!white, colframe=blue!50!black,
    title=Note, fonttitle=\bfseries, sharp corners, boxrule=0.5pt
}
\newtcolorbox{warnbox}{
    colback=orange!5!white, colframe=orange!70!black,
    title=Important, fonttitle=\bfseries, sharp corners, boxrule=0.5pt
}
\newtcolorbox{tipbox}{
    colback=green!5!white, colframe=green!50!black,
    title=Tip, fonttitle=\bfseries, sharp corners, boxrule=0.5pt
}

\titleformat{\section}{\Large\bfseries\color{blue!50!black}}{\thesection}{1em}{}
\titleformat{\subsection}{\large\bfseries\color{blue!40!black}}{\thesubsection}{1em}{}

\title{\textbf{Deploying ML Models on AWS ECS and EKS}\\
\large A Step-by-Step Container Orchestration Guide}
\author{Prabal}
\date{\today}

\begin{document}
\maketitle
\tableofcontents
\newpage

% =========================================================
\section{Introduction}
% =========================================================

When your ML workload outgrows SageMaker's managed abstractions or Lambda's
size and GPU limitations, AWS offers two container orchestrators for running
custom inference stacks at scale:

\begin{itemize}[leftmargin=*]
    \item \textbf{Amazon ECS} (Elastic Container Service) --- AWS-native
          orchestrator. Simpler mental model, tight integration with the AWS
          ecosystem, and lower operational overhead.
    \item \textbf{Amazon EKS} (Elastic Kubernetes Service) --- managed
          Kubernetes control plane. More portable, more powerful, and backed by
          the enormous Kubernetes ecosystem (Helm, Argo, KServe, Kubeflow),
          at the cost of higher complexity.
\end{itemize}

Both pull container images from \textbf{Amazon ECR} and can run on two
underlying compute platforms:

\begin{itemize}[leftmargin=*,nosep]
    \item \textbf{Fargate} --- serverless containers, no EC2 to manage. No GPU.
    \item \textbf{EC2} --- you manage the instance fleet; required for GPU,
          Inferentia, spot, and custom kernels.
\end{itemize}

\begin{notebox}
\textbf{When to choose ECS / EKS over SageMaker:}
\begin{itemize}[leftmargin=*,nosep]
    \item You need a custom serving stack --- Triton Inference Server,
          TorchServe, Ray Serve, vLLM, Text Generation Inference.
    \item You already run Kubernetes for the rest of your platform.
    \item You want fine-grained control over networking, sidecars, service mesh.
    \item Cost: EC2 reservations / spot can beat SageMaker list prices.
    \item You are serving LLMs or diffusion models and need exotic GPU
          topologies (multi-GPU, tensor parallelism).
\end{itemize}
\end{notebox}

% =========================================================
\section{ECS vs.\ EKS --- Choosing}
% =========================================================

\begin{longtable}{p{4cm} p{5cm} p{5.5cm}}
\toprule
\textbf{Dimension} & \textbf{ECS} & \textbf{EKS} \\
\midrule
Control plane cost & Free & \$0.10/hour/cluster ($\sim$\$73/mo) \\
\addlinespace
Learning curve & Low & High (Kubernetes) \\
\addlinespace
Portability & AWS only & Runs anywhere Kubernetes runs \\
\addlinespace
Ecosystem & AWS-centric & Helm, Argo, KServe, Kubeflow, Istio \\
\addlinespace
Autoscaling & Service + cluster autoscaling & HPA, VPA, KEDA, Karpenter \\
\addlinespace
Best for & Small/medium teams, AWS-only shops & Platform teams, multi-tenant, advanced ML stacks \\
\bottomrule
\end{longtable}

\begin{tipbox}
\textbf{Rule of thumb:} Start with ECS on Fargate. Move to EKS only when you
hit a specific need Kubernetes uniquely solves --- Helm-based OSS inference
stacks (KServe, Ray Serve, vLLM), multi-model GPU sharing, or an existing
platform engineering investment.
\end{tipbox}

% =========================================================
\section{Common Prerequisites}
% =========================================================

\begin{enumerate}[leftmargin=*]
    \item AWS CLI configured, Docker installed and running.
    \item A VPC with at least two subnets across different Availability Zones.
    \item An ECR repository holding your inference image (see
          \S\ref{sec:image}).
    \item A trained model artifact (baked in image \emph{or} stored in S3 / EFS /
          FSx for Lustre).
    \item IAM permissions to create ECS clusters / EKS clusters, task roles,
          load balancers, autoscaling policies.
\end{enumerate}

% =========================================================
\section{Building the Inference Image}
\label{sec:image}
% =========================================================

Unlike Lambda, there is no AWS-provided runtime wrapper --- your container
simply listens on an HTTP port and responds to health checks. A minimal
FastAPI + PyTorch image:

\begin{lstlisting}[style=dockercode]
FROM python:3.11-slim

WORKDIR /app

# System deps (if needed for your model)
RUN apt-get update && apt-get install -y --no-install-recommends \
    build-essential && rm -rf /var/lib/apt/lists/*

COPY requirements.txt .
RUN pip install --no-cache-dir -r requirements.txt

COPY app/ ./app/
COPY model/ ./model/

EXPOSE 8080
CMD ["uvicorn", "app.main:app", "--host", "0.0.0.0", "--port", "8080", \
     "--workers", "1"]
\end{lstlisting}

\textbf{Minimal FastAPI app (\texttt{app/main.py}):}

\begin{lstlisting}[style=pycode]
from fastapi import FastAPI, HTTPException
from pydantic import BaseModel
import torch

app = FastAPI()
MODEL = torch.load("model/model.pt", map_location="cpu")
MODEL.eval()

class Request(BaseModel):
    inputs: list

@app.get("/health")
def health():
    return {"status": "ok"}

@app.post("/predict")
def predict(req: Request):
    try:
        x = torch.tensor(req.inputs, dtype=torch.float32)
        with torch.inference_mode():
            y = MODEL(x)
        return {"predictions": y.tolist()}
    except Exception as e:
        raise HTTPException(status_code=500, detail=str(e))
\end{lstlisting}

\textbf{Push to ECR} (same steps as Lambda doc):

\begin{lstlisting}[style=bashcode]
export AWS_REGION=us-east-1
export ACCOUNT_ID=$(aws sts get-caller-identity --query Account --output text)
export REPO=ml-inference

aws ecr create-repository --repository-name $REPO --region $AWS_REGION

aws ecr get-login-password --region $AWS_REGION | \
  docker login --username AWS \
  --password-stdin $ACCOUNT_ID.dkr.ecr.$AWS_REGION.amazonaws.com

docker build --platform linux/amd64 -t $REPO:v1 .
docker tag  $REPO:v1  $ACCOUNT_ID.dkr.ecr.$AWS_REGION.amazonaws.com/$REPO:v1
docker push $ACCOUNT_ID.dkr.ecr.$AWS_REGION.amazonaws.com/$REPO:v1
\end{lstlisting}

\begin{tipbox}
For production-grade model servers, skip DIY FastAPI and use a purpose-built
server: \textbf{NVIDIA Triton} (multi-framework, dynamic batching, model
versioning), \textbf{TorchServe}, \textbf{TensorFlow Serving},
\textbf{vLLM} (LLMs), or \textbf{Text Generation Inference}. They handle batching,
multi-model hosting, and metrics far better than a handwritten server.
\end{tipbox}

% =========================================================
\section{Path A --- Deploy on ECS (Fargate)}
% =========================================================

ECS models a workload as:

\begin{itemize}[leftmargin=*,nosep]
    \item \textbf{Task Definition} --- the ``recipe'': image URI, CPU/memory,
          env vars, ports, logging.
    \item \textbf{Service} --- long-running controller that keeps N tasks
          healthy behind a load balancer.
    \item \textbf{Cluster} --- logical grouping of services and tasks.
\end{itemize}

% ---------------------------------
\subsection{Step A1 --- Create the ECS Cluster}
% ---------------------------------

\begin{lstlisting}[style=bashcode]
aws ecs create-cluster \
    --cluster-name ml-cluster \
    --capacity-providers FARGATE FARGATE_SPOT \
    --default-capacity-provider-strategy \
      capacityProvider=FARGATE,weight=1
\end{lstlisting}

% ---------------------------------
\subsection{Step A2 --- IAM Roles}
% ---------------------------------

You need two roles:

\begin{itemize}[leftmargin=*,nosep]
    \item \textbf{Task Execution Role} --- used by ECS itself to pull from ECR
          and write logs to CloudWatch. Attach
          \texttt{AmazonECSTaskExecutionRolePolicy}.
    \item \textbf{Task Role} --- assumed by the \emph{application} inside the
          container. Grant it whatever AWS permissions your inference code
          needs (e.g., \texttt{s3:GetObject} on the model bucket).
\end{itemize}

% ---------------------------------
\subsection{Step A3 --- Task Definition}
% ---------------------------------

\texttt{task-def.json}:

\begin{lstlisting}[style=jsoncode]
{
  "family": "ml-inference",
  "networkMode": "awsvpc",
  "requiresCompatibilities": ["FARGATE"],
  "cpu": "2048",
  "memory": "8192",
  "executionRoleArn": "arn:aws:iam::ACCOUNT:role/ecsTaskExecutionRole",
  "taskRoleArn":      "arn:aws:iam::ACCOUNT:role/ml-inference-task-role",
  "containerDefinitions": [{
    "name": "ml",
    "image": "ACCOUNT.dkr.ecr.us-east-1.amazonaws.com/ml-inference:v1",
    "essential": true,
    "portMappings": [{"containerPort": 8080, "protocol": "tcp"}],
    "healthCheck": {
      "command": ["CMD-SHELL", "curl -f http://localhost:8080/health || exit 1"],
      "interval": 30, "timeout": 5, "retries": 3, "startPeriod": 30
    },
    "logConfiguration": {
      "logDriver": "awslogs",
      "options": {
        "awslogs-group": "/ecs/ml-inference",
        "awslogs-region": "us-east-1",
        "awslogs-stream-prefix": "ecs"
      }
    }
  }]
}
\end{lstlisting}

Register it:

\begin{lstlisting}[style=bashcode]
aws ecs register-task-definition --cli-input-json file://task-def.json
\end{lstlisting}

\begin{warnbox}
Fargate CPU/memory combinations are \emph{not} arbitrary. Valid pairs include
0.25 vCPU / 0.5--2 GB, 2 vCPU / 4--16 GB, 16 vCPU / 32--120 GB, etc. Check
the Fargate task sizing table before picking values or the task will fail to
register.
\end{warnbox}

% ---------------------------------
\subsection{Step A4 --- Application Load Balancer}
% ---------------------------------

Create an internet-facing ALB in your public subnets, then a target group of
type \texttt{ip} (required for Fargate with \texttt{awsvpc} networking):

\begin{lstlisting}[style=bashcode]
aws elbv2 create-target-group \
    --name ml-tg \
    --protocol HTTP --port 8080 \
    --vpc-id vpc-XXXXXXXX \
    --target-type ip \
    --health-check-path /health \
    --health-check-interval-seconds 30
\end{lstlisting}

Attach a listener forwarding :80 (or :443 with an ACM cert) to the target group.

% ---------------------------------
\subsection{Step A5 --- Create the Service}
% ---------------------------------

\begin{lstlisting}[style=bashcode]
aws ecs create-service \
    --cluster ml-cluster \
    --service-name ml-inference \
    --task-definition ml-inference \
    --desired-count 2 \
    --launch-type FARGATE \
    --network-configuration "awsvpcConfiguration={
        subnets=[subnet-A,subnet-B],
        securityGroups=[sg-XXXX],
        assignPublicIp=DISABLED}" \
    --load-balancers "targetGroupArn=arn:...:targetgroup/ml-tg/...,\
        containerName=ml,containerPort=8080" \
    --deployment-configuration \
        "minimumHealthyPercent=100,maximumPercent=200"
\end{lstlisting}

ECS will launch two tasks, register them with the ALB, and replace any that
fail health checks.

% ---------------------------------
\subsection{Step A6 --- Service Autoscaling}
% ---------------------------------

Register the service as a scalable target and attach a target-tracking policy:

\begin{lstlisting}[style=bashcode]
aws application-autoscaling register-scalable-target \
    --service-namespace ecs \
    --resource-id service/ml-cluster/ml-inference \
    --scalable-dimension ecs:service:DesiredCount \
    --min-capacity 2 --max-capacity 20

aws application-autoscaling put-scaling-policy \
    --policy-name ml-cpu-scale \
    --service-namespace ecs \
    --resource-id service/ml-cluster/ml-inference \
    --scalable-dimension ecs:service:DesiredCount \
    --policy-type TargetTrackingScaling \
    --target-tracking-scaling-policy-configuration file://scale.json
\end{lstlisting}

Where \texttt{scale.json} targets 60\% CPU utilization with cooldowns.

\begin{tipbox}
For latency-sensitive services, scale on
\texttt{ALBRequestCountPerTarget} or a custom CloudWatch metric
(e.g., queue depth) rather than CPU --- CPU reacts too slowly and often stays
flat until the service is already saturated.
\end{tipbox}

% ---------------------------------
\subsection{Step A7 --- Rolling Updates}
% ---------------------------------

Push a new image tag, register a new task-definition revision, then:

\begin{lstlisting}[style=bashcode]
aws ecs update-service \
    --cluster ml-cluster \
    --service ml-inference \
    --task-definition ml-inference:2 \
    --force-new-deployment
\end{lstlisting}

With \texttt{minimumHealthyPercent=100, maximumPercent=200}, ECS launches new
tasks alongside old ones and drains the old ones only after new ones are
healthy (blue/green style).

For canary / linear / all-at-once with automatic rollback, switch the
deployment controller to \texttt{CODE\_DEPLOY} and configure CodeDeploy.

% =========================================================
\section{Path B --- Deploy on EKS}
% =========================================================

EKS is managed Kubernetes. The ML-relevant objects are:

\begin{itemize}[leftmargin=*,nosep]
    \item \textbf{Deployment} --- declarative pod template + replica count.
    \item \textbf{Service} --- stable in-cluster IP / DNS.
    \item \textbf{Ingress} / \textbf{ALB Ingress} --- external HTTPS entry point.
    \item \textbf{HPA} --- Horizontal Pod Autoscaler, scales pods.
    \item \textbf{Karpenter} / \textbf{Cluster Autoscaler} --- scales nodes.
\end{itemize}

% ---------------------------------
\subsection{Step B1 --- Create the EKS Cluster}
% ---------------------------------

The easiest path is \texttt{eksctl}:

\begin{lstlisting}[style=bashcode]
eksctl create cluster \
    --name ml-eks \
    --region us-east-1 \
    --version 1.29 \
    --nodegroup-name gpu-nodes \
    --node-type g5.xlarge \
    --nodes 2 --nodes-min 1 --nodes-max 10 \
    --with-oidc \
    --ssh-access --ssh-public-key my-key
\end{lstlisting}

For Fargate-only (CPU, no GPU):

\begin{lstlisting}[style=bashcode]
eksctl create cluster \
    --name ml-eks \
    --fargate \
    --region us-east-1
\end{lstlisting}

\begin{warnbox}
The EKS control plane costs \textbf{\$0.10/hour} ($\sim$\$73/month) \emph{per
cluster}, on top of node costs. Do not spin up clusters casually for dev.
Share clusters across teams where possible.
\end{warnbox}

% ---------------------------------
\subsection{Step B2 --- GPU Support}
% ---------------------------------

For GPU nodes (e.g., \texttt{g5}, \texttt{p4d}, \texttt{p5}), install the
NVIDIA device plugin DaemonSet:

\begin{lstlisting}[style=bashcode]
kubectl apply -f \
  https://raw.githubusercontent.com/NVIDIA/k8s-device-plugin/v0.14.0/nvidia-device-plugin.yml
\end{lstlisting}

Pods then request GPUs via \texttt{resources.limits["nvidia.com/gpu"]: 1}.

% ---------------------------------
\subsection{Step B3 --- IAM Roles for Service Accounts (IRSA)}
% ---------------------------------

IRSA maps a Kubernetes ServiceAccount to an IAM role, so only pods running as
that SA get those AWS permissions --- the least-privilege pattern on EKS.

\begin{lstlisting}[style=bashcode]
eksctl create iamserviceaccount \
    --cluster ml-eks \
    --name ml-inference-sa \
    --namespace ml \
    --attach-policy-arn arn:aws:iam::aws:policy/AmazonS3ReadOnlyAccess \
    --approve
\end{lstlisting}

% ---------------------------------
\subsection{Step B4 --- Deployment Manifest}
% ---------------------------------

\texttt{deployment.yaml}:

\begin{lstlisting}[style=yamlcode]
apiVersion: apps/v1
kind: Deployment
metadata:
  name: ml-inference
  namespace: ml
spec:
  replicas: 2
  strategy:
    type: RollingUpdate
    rollingUpdate:
      maxSurge: 1
      maxUnavailable: 0
  selector:
    matchLabels:
      app: ml-inference
  template:
    metadata:
      labels:
        app: ml-inference
    spec:
      serviceAccountName: ml-inference-sa
      containers:
      - name: ml
        image: ACCOUNT.dkr.ecr.us-east-1.amazonaws.com/ml-inference:v1
        ports:
        - containerPort: 8080
        resources:
          requests:
            cpu: "1"
            memory: "4Gi"
            nvidia.com/gpu: "1"    # remove for CPU-only
          limits:
            cpu: "2"
            memory: "8Gi"
            nvidia.com/gpu: "1"
        readinessProbe:
          httpGet: {path: /health, port: 8080}
          initialDelaySeconds: 30
          periodSeconds: 10
        livenessProbe:
          httpGet: {path: /health, port: 8080}
          initialDelaySeconds: 60
          periodSeconds: 30
\end{lstlisting}

% ---------------------------------
\subsection{Step B5 --- Service and Ingress}
% ---------------------------------

\texttt{service.yaml}:

\begin{lstlisting}[style=yamlcode]
apiVersion: v1
kind: Service
metadata:
  name: ml-inference
  namespace: ml
spec:
  type: ClusterIP
  selector:
    app: ml-inference
  ports:
  - port: 80
    targetPort: 8080
    protocol: TCP
\end{lstlisting}

External exposure via the \textbf{AWS Load Balancer Controller} (install via
Helm once per cluster), which provisions a real AWS ALB from an Ingress:

\begin{lstlisting}[style=yamlcode]
apiVersion: networking.k8s.io/v1
kind: Ingress
metadata:
  name: ml-inference
  namespace: ml
  annotations:
    kubernetes.io/ingress.class: alb
    alb.ingress.kubernetes.io/scheme: internet-facing
    alb.ingress.kubernetes.io/target-type: ip
    alb.ingress.kubernetes.io/healthcheck-path: /health
spec:
  rules:
  - http:
      paths:
      - path: /
        pathType: Prefix
        backend:
          service:
            name: ml-inference
            port:
              number: 80
\end{lstlisting}

Apply everything:

\begin{lstlisting}[style=bashcode]
kubectl create namespace ml
kubectl apply -f deployment.yaml
kubectl apply -f service.yaml
kubectl apply -f ingress.yaml
\end{lstlisting}

% ---------------------------------
\subsection{Step B6 --- Horizontal Pod Autoscaler}
% ---------------------------------

\begin{lstlisting}[style=yamlcode]
apiVersion: autoscaling/v2
kind: HorizontalPodAutoscaler
metadata:
  name: ml-inference
  namespace: ml
spec:
  scaleTargetRef:
    apiVersion: apps/v1
    kind: Deployment
    name: ml-inference
  minReplicas: 2
  maxReplicas: 20
  metrics:
  - type: Resource
    resource:
      name: cpu
      target:
        type: Utilization
        averageUtilization: 60
\end{lstlisting}

For queue-depth-based scaling, custom metrics, or scale-to-zero, install
\textbf{KEDA} --- it supports SQS, Kafka, Prometheus, CloudWatch, and many
other external triggers.

% ---------------------------------
\subsection{Step B7 --- Node Autoscaling}
% ---------------------------------

Two options:

\begin{itemize}[leftmargin=*]
    \item \textbf{Cluster Autoscaler} --- reactive, per-node-group, simple.
    \item \textbf{Karpenter} --- modern AWS-native autoscaler; picks the
          cheapest suitable instance type dynamically (great for mixed
          CPU/GPU fleets and spot). Recommended for new clusters.
\end{itemize}

% ---------------------------------
\subsection{Step B8 --- Purpose-Built ML Serving on EKS}
% ---------------------------------

Rather than writing your own Deployments, many teams deploy pre-built
inference platforms on EKS:

\begin{itemize}[leftmargin=*]
    \item \textbf{KServe} --- standard CRDs (\texttt{InferenceService}) with
          canary, autoscaling, explainability, batching. The de-facto model
          serving layer on Kubernetes.
    \item \textbf{Ray Serve} --- Python-first, great for composition and
          multi-model pipelines.
    \item \textbf{Seldon Core} --- enterprise-grade inference graphs.
    \item \textbf{NVIDIA Triton} on EKS --- best-in-class GPU utilization,
          dynamic batching, multi-framework.
    \item \textbf{vLLM / TGI} --- LLM-specific servers with paged attention
          and continuous batching.
\end{itemize}

% =========================================================
\section{GPU and Accelerator Notes}
% =========================================================

\begin{longtable}{p{3.5cm} p{5cm} p{6cm}}
\toprule
\textbf{Instance} & \textbf{Accelerator} & \textbf{Use case} \\
\midrule
\texttt{g4dn.*} & NVIDIA T4 (16 GB) & Low-cost inference, small-to-medium DL \\
\texttt{g5.*}   & NVIDIA A10G (24 GB) & Mid-sized LLMs ($\leq$13B), SD, CV \\
\texttt{p4d.*}  & 8$\times$ A100 (40 GB) & Large LLMs, training, tensor parallel \\
\texttt{p5.*}   & 8$\times$ H100 (80 GB) & Frontier models, long context \\
\texttt{inf1.*} & AWS Inferentia 1 & Cheapest \$/inference for supported models \\
\texttt{inf2.*} & AWS Inferentia 2 & Larger models, generative AI \\
\texttt{trn1.*} & AWS Trainium & Training (occasionally inference) \\
\bottomrule
\end{longtable}

\begin{tipbox}
For steady-state LLM serving, \texttt{inf2} with the Neuron SDK is often
30--50\% cheaper than equivalent NVIDIA GPUs --- at the cost of a narrower
model-support matrix and a model-compilation step.
\end{tipbox}

% =========================================================
\section{Model Weight Storage Strategies}
% =========================================================

\begin{enumerate}[leftmargin=*]
    \item \textbf{Baked in image} --- simplest, deterministic. Image rebuilds
          on every model update. Good for $\leq$ few GB.
    \item \textbf{S3 at startup} --- image stays small; cold start downloads
          weights to local disk. Typical for 5--50 GB models.
    \item \textbf{EFS mount} --- shared read-only filesystem across pods. Great
          for multi-model endpoints. Medium throughput.
    \item \textbf{FSx for Lustre} --- high-throughput parallel filesystem
          backed by S3. Standard for large-scale GPU inference (LLMs) where
          you need many GB loaded fast across many pods.
\end{enumerate}

% =========================================================
\section{Observability}
% =========================================================

\paragraph{ECS.} CloudWatch Container Insights (one flag at cluster creation),
plus standard CloudWatch Logs from the \texttt{awslogs} driver. AWS
Distro for OpenTelemetry (ADOT) for traces and custom metrics.

\paragraph{EKS.} Install the standard Kubernetes observability stack:
\begin{itemize}[leftmargin=*,nosep]
    \item \textbf{Prometheus} (AMP --- Amazon Managed Prometheus) + \textbf{Grafana}
          (AMG) for metrics.
    \item \textbf{Fluent Bit} $\rightarrow$ CloudWatch Logs / OpenSearch.
    \item \textbf{AWS Distro for OpenTelemetry} for traces.
    \item \textbf{DCGM Exporter} for GPU metrics (utilization, memory,
          temperature, throttling).
\end{itemize}

% =========================================================
\section{Cost Considerations}
% =========================================================

\begin{itemize}[leftmargin=*]
    \item \textbf{Spot instances} --- up to 90\% off on-demand; EKS + Karpenter
          handles interruption gracefully. Great for batch or stateless
          inference replicas.
    \item \textbf{Savings Plans} for steady 24/7 workloads (1-year or 3-year
          commitment, 30--50\% savings).
    \item \textbf{Right-size continuously} --- most teams over-provision memory
          by 2--3$\times$. Use Container Insights or Kubecost to find slack.
    \item \textbf{Scale to zero} overnight / off-hours via KEDA cron scalers or
          scheduled ECS service updates.
    \item \textbf{Share GPUs} via MIG (A100/H100) or NVIDIA Time-Slicing for
          small models --- up to 7$\times$ density on a single card.
\end{itemize}

% =========================================================
\section{Production Best Practices}
% =========================================================

\begin{enumerate}[leftmargin=*]
    \item \textbf{IaC everything} --- Terraform, CDK, Pulumi. Never click-ops
          a production cluster.
    \item \textbf{GitOps on EKS} --- Argo CD or Flux for declarative,
          auditable deployments from a Git repo.
    \item \textbf{Pod Disruption Budgets} on EKS so voluntary disruptions
          (node upgrades) don't take down all replicas.
    \item \textbf{Multi-AZ} --- spread replicas across $\geq$ 2 AZs with
          topology spread constraints / ECS placement strategies.
    \item \textbf{Warm replicas always $\geq$ 2} --- avoid single-point-of-failure
          on rolling updates or AZ failures.
    \item \textbf{Image scanning} --- enable ECR scan-on-push; block critical
          CVEs at CI time.
    \item \textbf{Dynamic batching} --- if using Triton or a custom server,
          enable micro-batching; doubles or triples throughput on GPU.
    \item \textbf{Separate training and inference clusters} --- different
          scaling profiles, security posture, and cost lifecycles.
\end{enumerate}

% =========================================================
\section{Troubleshooting Checklist}
% =========================================================

\begin{longtable}{p{5.5cm} p{9.2cm}}
\toprule
\textbf{Symptom} & \textbf{Likely Cause / Fix} \\
\midrule
ECS task stuck in \texttt{PENDING} &
Usually networking: the task cannot pull the image. Check that private
subnets have a NAT Gateway \emph{or} VPC endpoints for ECR + S3 + CloudWatch
Logs. \\
\addlinespace
\texttt{CannotPullContainerError} &
Task execution role missing ECR permissions, \emph{or} image is for the wrong
architecture. Rebuild with \texttt{--platform linux/amd64}. \\
\addlinespace
ALB targets unhealthy &
Health check path returns non-2xx, or pod takes longer than the initial delay
to start. Increase \texttt{initialDelaySeconds} / health check grace period. \\
\addlinespace
EKS pod pending, \texttt{Insufficient nvidia.com/gpu} &
No GPU nodes available or NVIDIA device plugin not installed. Check Karpenter
/ autoscaler node provisioning logs. \\
\addlinespace
OOMKilled in Kubernetes &
Memory request/limit too low --- framework lazy-allocates. Set \texttt{limits.memory}
generously or the kernel will reap the pod. \\
\addlinespace
GPU utilization near 0\% &
Tiny batch sizes; enable dynamic batching in Triton / TGI, or raise client
concurrency. \\
\addlinespace
Image pull is slow on scale-up &
Large image ($\geq$ 5 GB). Use ECR \emph{pull-through cache} + warm AMIs, or
mount weights from EFS/FSx and keep the code image small. \\
\bottomrule
\end{longtable}

% =========================================================
\section{Summary}
% =========================================================

\paragraph{ECS flow:}
\begin{enumerate}[leftmargin=*,nosep]
    \item Build and push image to ECR.
    \item Create ECS cluster.
    \item Define task execution and task roles.
    \item Register task definition (image, CPU/memory, ports, logging).
    \item Create ALB + target group.
    \item Create ECS Service with desired count, networking, load balancer.
    \item Attach Application Auto Scaling policy.
    \item Update via new task-definition revisions.
\end{enumerate}

\paragraph{EKS flow:}
\begin{enumerate}[leftmargin=*,nosep]
    \item Create cluster with \texttt{eksctl} (managed / Fargate / GPU node groups).
    \item Install AWS Load Balancer Controller, NVIDIA device plugin,
          Karpenter, metrics-server.
    \item Build and push image to ECR.
    \item Define Deployment, Service, Ingress manifests.
    \item Use IRSA for pod-level AWS permissions.
    \item Attach HPA (or KEDA) for pod scaling; Karpenter for node scaling.
    \item Optionally layer on KServe / Triton / Ray Serve for richer ML serving.
\end{enumerate}

\paragraph{The decision.} Choose \textbf{ECS} for simplicity and AWS-native
workflows; choose \textbf{EKS} when you need Kubernetes' ecosystem, portability,
or advanced serving primitives. In both cases, ECR is the image registry,
Fargate is the no-ops compute option, and EC2 (GPU / Inferentia) is the
high-performance compute option.

\end{document}
